\subsection{真实拍摄 VLM 探测实验}
\label{sec:vlm}

为检验传统 OCR 的结论能否推广至更强的识别器，我们使用了通过 SiliconFlow 应用程序编程接口（API）提供的三个商业视觉语言模型（VLM）端点：Qwen3-VL-32B-Instruct~\cite{b_qwen3vl2025}、Kimi-K2.6 和 GLM-4.5V~\cite{b_glm45v2025}。本节是一项\emph{边界探测}：若 VLM 能够恢复传统 OCR 无法识别的文本，则本方法的实际适用范围仅限于缓解传统 OCR 对短曝光图像的恢复。

探测实验覆盖三个距离下的 0$^{\circ}$ 正轴拍摄。0.5\,m 实验包含 156 个输入：36 个未保护短曝光基线、各 36 个高抑制短曝光和长曝光输入，以及从 60\,fps 视频生成的四种聚合视图（每种 12 个片段）。1.0\,m 和 1.5\,m 实验各包含 120 个高抑制输入：各 36 个短曝光和长曝光输入，外加四类视频聚合视图（每类 12 个片段）。VLM 输入与表~\ref{tab:real_vlm} 中 OCR 的最优引擎（best-of-engine，BoE）参照使用同一批文件，其标称标签为 3.91/31.25\,ms。所有 VLM 调用均使用相机原始 JPEG。\S\ref{sec:real_capture} 中固定的三引擎网格评估的是传统 OCR，而不是 VLM 输入预处理；先导实验中的 CLAHE 增强恢复出了更多 VLM 文本，因此基于原始 JPEG 的 VLM 结果并非攻击上界。

在 1,188 次模型请求中，1,133 次返回了可解析结果：0.5\,m 处失败 34 次，1.0\,m 处失败 21 次，1.5\,m 处无失败。每个表格单元均报告有效 $N$，且结果以响应可解析为条件。GLM 的失败集中在混合内容、数字和凭据条目上，因此不能将缺失视为随机缺失。为此，我们分别将每次失败调用的恢复率赋为 0 或 100\%，计算逐单元格的下界和上界。对于 GLM 在 0.5\,m 的短曝光单元格，字符恢复率的界限为 61.0--74.9\%，完全恢复率的界限为 44.4--58.3\%。我们不使用条件单元格对模型进行排序；所有单元格的完整界限均随分析一并归档。

我们将无失败的 1.5\,m 实验作为主要的跨模型证据；1.0\,m 实验则作为短曝光距离扫描的补充数据。在 1.0\,m 处，Qwen 和 Kimi 的长曝光完全匹配率分别为 84.4\% 和 74.3\%，三个模型的成功调用在 150 帧视频均值下得到 70.0--83.3\% 的完全匹配率。GLM 在 1.0\,m 处的长曝光完全匹配率仅为 3.3\%，而在 1.5\,m 处为 58.3\%；但依赖内容的失败使我们无法主张结果具有跨距离一致性。为保持紧凑，表~\ref{tab:real_vlm} 在 0.5\,m 处仅报告 150 帧视频均值；1.5\,m 实验则报告全部四种聚合视图，以展示结果对视图选择的敏感性。

所有调用均使用 SiliconFlow 的 OpenAI 兼容 API。记录的提供方标识符为 \texttt{Qwen/\allowbreak Qwen3-\allowbreak VL-\allowbreak 32B-\allowbreak Instruct}、\texttt{Pro/\allowbreak moonshotai/\allowbreak Kimi-\allowbreak K2.6} 和 \texttt{zai-\allowbreak org/\allowbreak GLM-\allowbreak 4.5V}。Qwen3-VL 和 GLM-4.5V 有公开技术报告~\cite{b_qwen3vl2025,b_glm45v2025}。对于 Kimi-K2.6，我们以提供方端点标识符和归档的请求元数据作为可复现性记录，而不对另有文档说明的上游公开版本作出主张。第二个标识符中的 \texttt{Pro/} 前缀为 SiliconFlow 所特有，该端点在我们归档的调用中接受图像输入。解码温度为 0，最大输出为 4,096 个词元。共用提示词要求仅输出可见文本并禁止猜测；模型不会获得真实值。完整提示词、参数和逐次调用输出均已归档。指标遵循 \S\ref{sec:setup}，并在拍摄样本层面汇总。公开的 CET6 文档可能因训练数据记忆而抬高恢复率，因此合成片段与完全匹配率能提供更干净的区分性证据。

\textbf{0.5\,m 会话选择透明性}：归档保留了同一位置的另一个会话（\path{real_capture_vlm_d0.5_a0_rerun.json}，时间戳 012715），其中高抑制短曝光输入几乎全黑，因而三个模型均给出空转录，字符恢复率和完全恢复率均为 0\%。表~\ref{tab:real_vlm} 使用 141107（3.91\,ms）会话，因为它与 OCR 参照使用同一批文件，且其输入并非全黑。这是仅有的两个 0.5\,m 正轴 VLM 会话，因此 Qwen 观测到的短曝光完全匹配结果从欠曝光会话的 0/36 个拍摄样本到表中会话的 28/36 个拍摄样本不等，并非稳定的单会话估计。充足曝光准则是在检查两个会话后才采用的，因此属于事后选择规则；保留两个会话并报告结果范围，可以明确揭示其曝光敏感性。表中会话是偏向攻击者的保守情形，另一个会话则保留作为欠曝光诊断。

\begin{table*}[!t]
\caption{真实拍摄 VLM 探测结果（0$^{\circ}$ 正轴；除原始短曝光基线外，所有行均为高抑制档）。每个 VLM 单元格报告完全匹配计数 / 字符恢复率 / $N$（恢复数 / \% / 有效调用数）；存在失败的单元格以响应可解析为条件。OCR BoE 给出同批次三引擎最优引擎参照（完全匹配率 / 字符恢复率，\%）。时域均值行对 150 个短曝光帧取平均；其他聚合视图在正文中定义。0.5\,m 高抑制短曝光行采用曝光充足且偏向攻击者的 141107 会话；仅有的另一个 0.5\,m 会话存在欠曝光，Qwen 的完全匹配结果为 0/36。}
\label{tab:real_vlm}
\centering
\footnotesize
\setlength{\tabcolsep}{4.2pt}
\begin{tabular}{llcccc}
\toprule
距离 & 条件 & OCR BoE & Qwen3-VL-32B & Kimi-K2.6 & GLM-4.5V \\
\midrule
\multicolumn{6}{l}{\textit{短曝光距离扫描}}\\
\midrule
0.5\,m & 原始短曝光 & 58.3~/~92.4 & 30~/~94.5~/~36 & 30~/~95.6~/~36 & 19~/~92.0~/~26 \\
0.5\,m & 短曝光 & 0.0~/~8.4 & 28~/~91.5~/~36 & 17~/~84.7~/~36 & 16~/~70.8~/~31 \\
1.0\,m & 短曝光 & 0.0~/~3.6 & 20~/~79.1~/~36 & 14~/~68.6~/~36 & 9~/~47.0~/~33 \\
1.5\,m & 短曝光 & 0.0~/~6.5 & 17~/~68.9~/~36 & 12~/~61.5~/~36 & 13~/~70.0~/~36 \\
\midrule
\multicolumn{6}{l}{\textit{关键攻击视图}}\\
\midrule
0.5\,m & 长曝光 & 0.0~/~2.1 & 28~/~85.7~/~34 & 1~/~7.9~/~36 & 0~/~4.9~/~30 \\
0.5\,m & 视频（150 帧均值） & 8.3~/~50.7 & 10~/~94.0~/~12 & 10~/~95.5~/~12 & 7~/~91.8~/~8 \\
1.5\,m & 长曝光 & 0.0~/~13.2 & 33~/~89.3~/~36 & 32~/~89.9~/~36 & 21~/~62.4~/~36 \\
1.5\,m & 视频（单帧最优） & 0.0~/~4.3 & 1~/~24.8~/~12 & 0~/~12.7~/~12 & 0~/~29.8~/~12 \\
1.5\,m & 视频（150 帧均值） & 0.0~/~48.4 & 8~/~90.3~/~12 & 8~/~90.3~/~12 & 7~/~88.5~/~12 \\
1.5\,m & 视频（窗口均值最优） & 0.0~/~8.6 & 7~/~89.0~/~12 & 5~/~84.2~/~12 & 8~/~87.2~/~12 \\
1.5\,m & 视频（最大值投影） & 0.0~/~19.0 & 8~/~90.1~/~12 & 7~/~89.2~/~12 & 7~/~88.7~/~12 \\
\bottomrule
\end{tabular}
\end{table*}

\begin{table}[!t]
\caption{0.5\,m 高抑制短曝光条件下按内容类型划分的 VLM 字符恢复率。基线列给出同一会话中 Qwen 对未保护原始短曝光图像的结果（Kimi/GLM 与之相近）；GLM 偶有调用失败，括号内的 $N$ 表示有效数量。在 1.5\,m 处，三个模型对 CET6 密集页面的字符恢复率均为 0\%。}
\label{tab:vlm_content}
\centering
\footnotesize
\setlength{\tabcolsep}{2.8pt}
\begin{tabular}{lcccc}
\toprule
内容类型（$N$） & 基线（\%） & Qwen（\%） & Kimi（\%） & GLM（\%） \\
\midrule
CET6 文档（3） & 64.4 & 18.5 & 5.6 & 0.4 \\
凭据（6） & 97.6 & 97.6 & 88.3 & 77.1~(5) \\
代码片段（6） & 89.8 & 94.9 & 95.8 & 88.9 \\
数字字符串（6） & 100.0 & 100.0 & 98.6 & 50.0~(4) \\
混合内容（6） & 100.0 & 100.0 & 92.5 & 98.2~(4) \\
段落（6） & 100.0 & 99.8 & 81.7 & 66.5 \\
英文句子（3） & 95.3 & 95.3 & 96.9 & 94.6 \\
\bottomrule
\end{tabular}
\end{table}

表~\ref{tab:real_vlm} 给出了汇总的 VLM 探测结果：短曝光行展示距离变化，长曝光和视频聚合行展示积分攻击边界。表~\ref{tab:vlm_content} 按内容类型细分了 0.5\,m 高抑制短曝光条件下的字符恢复率。

\textbf{主要发现}：

\begin{enumerate}
\item \textbf{VLM 限定而非扩展了主张范围}：对于同一批 0.5\,m 照片，在高抑制短曝光条件下，最优引擎 OCR 的字符恢复率仅为 8.4\%，完全匹配率为 0\%；而 Qwen3-VL 的字符恢复率达到 91.5\%，并在来自 12 个内容项的 36 个拍摄样本中完全恢复 28 个。在 1.5\,m 处，三个模型仍保持 33.3--47.2\% 的完全匹配率和 61.5--70.0\% 的字符恢复率。即使提示词禁止猜测，VLM 仍可从可见片段中恢复大量文本。因此，在这项正轴、高抑制、三距离探测实验中，针对传统 OCR 的主张不能推广至 VLM 攻击者；本实验未评估离轴 VLM 拍摄和其他保护档位。

\item \textbf{VLM 攻击下的保护失效集中在大字号短片段}：在 0.5\,m 处，VLM 对 CET6 文档页面的字符恢复率为 0.4--18.5\%，而对凭据、代码和数字字符串的字符恢复率为 50--100\%（表~\ref{tab:vlm_content}）。在 1.5\,m 处，三个模型对 CET6 页面的恢复率均降至 0\%，但数字字符串和短句仍可恢复。这与字体尺度和语言先验的影响一致~\cite{b_shi2023,b_jiang2025}，但本实验并未单独识别其中任何一种机制。

\item \textbf{攻击能力随模型和条件而变化}：在无失败的 1.5\,m 会话中，标称长曝光完全恢复率为 58.3\%--91.7\%，字符恢复率为 62.4--89.9\%。这一差异仅为描述性结果，不应解释为稳定的模型排名。VLM 行为取决于端点版本、距离、拍摄视图和内容。

\item \textbf{持续视频聚合构成 VLM 失效边界}：在 0.5\,m 处，三个模型的 150 帧均值缺失界限覆盖 58.3--91.7\% 的完全匹配率和 61.2--95.5\% 的字符恢复率。所有聚合视图均来自同一段连拍，不能代表相互独立的短时攻击。在无失败的 1.5\,m 会话中，150 帧均值产生 58.3--66.7\% 的完全匹配率和 88.5--90.3\% 的字符恢复率，而单帧最优视图的完全恢复率为 0--8.3\%。本实验未确定所需的最短录制时长。
\end{enumerate}

因此，有效性主张必须由内容和攻击者类型共同限定：在本探测实验中，高抑制档降低了传统 OCR 对密集小字号文本的恢复率；但当攻击者采用现代 VLM 时，大字号短片段（凭据、代码、数字字符串）仍可恢复。这一负面结果将方法的实际适用范围从“阻止任意机器读取屏幕”重新限定为“降低传统 OCR 对短曝光批量图像的恢复”，并将面向 VLM 的显示保护明确为开放问题。对策方向见 \S\ref{sec:vlm_defense}。本节仅覆盖单摄像头、0$^{\circ}$ 正轴拍摄、一个档位、三个商业 API 模型，且每个单元格仅有 12 个视频片段；这些数值是描述性的边界探测结果，而非完整的跨几何条件或跨模型版本评估。

% Defined here for source locality with the real-capture protocol, but emitted
% after all attack experiments so the evaluation narrative remains continuous.
\long\def\deferreduserstudy{
% ----------------------------------------------------------
\subsection{用户体验研究}
\label{sec:user_study}

为评估受保护显示对人类可读性和视觉舒适度的影响，我们设计了一项基于 Web 的用户研究。

\textbf{设备与环境协议}：采用与 \S\ref{sec:setup} 相同的 240\,Hz 显示器。正式研究版本将实测刷新率的硬性门槛设为 $\geq$200\,Hz，以确保固定 $n=4$ 时基本子帧循环频率满足 $f_r/n\geq50$\,Hz；刷新率低于该门槛时将阻止参与，而不是在不提示的情况下减小 $n$。144--199\,Hz 的自适应路径仅用于明确标注的演示模式，其会话默认不纳入分析。每位参与者开始实验前，实验人员确认显示器亮度固定，自动亮度/省电模式和可变刷新率已禁用，浏览器统一为全屏模式，高负载后台进程已关闭，观看距离约为 60\,cm。显示模式发生任何变化后，均重新测量刷新率。

浏览器在每次试验的 \path{requestAnimationFrame} 循环中记录实际帧间隔；若相邻回调的间隔超过根据试验前实测刷新率计算的预期间隔的 1.5$\times$，则记录一次丢帧事件，同时记录估计丢帧数、丢帧率、平均/最大帧间隔、观测刷新率、基本子帧循环频率，以及包含反相帧的完整循环频率。这些记录仅反映浏览器的呈现节奏，不能替代对面板扫描或光度学特性的测量。

\textbf{伦理与安全}：拟开展的研究属于最小风险的人体参与研究，目前尚未收集任何参与者数据。招募开始前，作者将从具备权限的机构或院系主管部门取得并记录书面批准、豁免或免审认定；在取得该认定前，不会开始正式数据采集。参与者必须分别勾选知情同意与光敏感筛查：页面明确说明掩蔽显示涉及快速时序调制，若出现不适、眼疲劳、眩晕、恶心或头痛，应立即停止，并可随时无条件退出。自述有光敏性癫痫史或对闪烁敏感者不得继续参与。系统记录同意时间戳和筛查通过标记。实现采用保守的 $f_r/n\geq50$\,Hz 运行下限，但这一工程门槛不构成伦理或临床安全认定。在 240\,Hz 面板上取 $n=4$ 时，基本循环频率为 60\,Hz；反相帧将完整循环频率降至 $f_r/(n+1)=48$\,Hz，实现会单独记录警告和实际时序。注册信息（学号、姓名）仅用于参与管理；分析和发表仅使用去标识化数据。年龄和性别为可选的人口统计字段。

\textbf{任务设计}：每位参与者首先完成一次 12 秒、不计分、无掩蔽的 Canvas 打字热身；随后观看 10 秒、不计分的可读性优先完整掩蔽预览；再在同一可读性优先掩蔽下完成一次 8 秒、不计分的打字练习。这一流程将初始刺激适应和输入适应与随后四次各 20 秒的计分试验（被试内设计）分离开来。在每个打字页面上，包括热身、练习和计分试验，参与者点击“开始”前，刺激 Canvas 已可见且播放器已启动。点击“开始”后触发 5 秒倒计时，此时输入仍被禁用；仅在倒计时结束后，计时打字阶段和输入功能才同时开始。控制条件和掩蔽条件均采用相同的开始前观看流程。

两个条件各重复两次。实验人员连续分配的从零开始的序列号 $k$ 用于选择 ABBA 或 BAAB 顺序，以平衡练习效应与疲劳效应：(A)~\textbf{原始文本}（控制条件），以 $n=1$ 的无掩蔽 Canvas 显示；(B)~\textbf{掩蔽文本}，使用可读性优先配置：$n=4$、掩蔽加噪声、强抗 OCR（$\alpha_s=0.10$、$\alpha_g=0.12$、6 个循环）+ 弱反相（$\alpha=0.2$）。打字顺序索引为 $k\bmod2$，主观评分的拉丁方行索引为 $\lfloor k/2\rfloor\bmod6$。因此，每连续 12 名参与者恰好覆盖全部 $2\times6$ 种联合分配。正式会话对 $k$ 强制执行唯一性约束，后端会重新计算并核验两个顺序索引。

两个条件均由同一 \texttt{render\allowbreak Source\allowbreak Canvas} 路径渲染，使用相同的深色背景、浅色文字、Canvas 尺寸、23\,px 粗体字体和颜色。唯一的处理差异是时序保护流水线，从而避免极性和字重混杂。此处的“循环”是指 Web 实现依次轮换六组不同的掩蔽、噪声和条纹实例，随后再重复。这一 Web 播放特有设置与 \S\ref{sec:method} 中的真实拍摄配置共用 $\alpha_s$ 和 $\alpha_g$，并避免反复积分同一固定叠加图案。

评分不再逐位置比较字符，而是使用 Levenshtein 最小字符串距离（minimum string distance，MSD）~\cite{b_soukoreff2003}，将完整转录与最优目标前缀对齐；未输入的目标后缀不计为错误。对齐后得到准确率、每分钟字符数（CPM）和每分钟词数（WPM），同时记录编辑距离、对齐前缀长度和评分版本。首次按键延迟定义为从 5 秒倒计时结束时（此时输入启用且计分计时器启动）到首次使转录内容非空的输入事件之间的时间。分析首先在每名参与者内部对每个条件下的两次试验取平均，随后进行被试内配对比较。

\textbf{主观评分}：打字试验结束后，参与者针对界面显示的三个 1--5 级 Likert 条目，对 6~种条件进行评分：可读性（1 = 难以阅读，5 = 非常清晰）、稳定性（存储在 \texttt{flicker} 字段中；1 = 强烈闪烁，5 = 无可感知闪烁）和即时视觉舒适度（\texttt{fatigue}；1 = 非常不适，5 = 非常舒适）。因此，每个条目的分值越高均代表结果越好。条件包括无掩蔽原始锚点、$n\in\{2,3,4\}$ 的掩蔽加噪声、$n=4$ 的仅掩蔽，以及打字任务所用的同一 $n=4$ 可读性优先完整配置（掩蔽加噪声+强抗 OCR+弱反相）。前三种掩蔽层级在 240\,Hz 面板上均能达到 $\geq$60\,Hz 的真实时序播放；原计划的 $n=8$ 只会产生 30\,Hz 并触发静态回退，因此不纳入时序消融。六种条件按 6 阶平衡拉丁方呈现；每种条件至少观看 10~秒后方可提交，同时记录观看开始/结束时间和持续时长。无掩蔽锚点也用于辅助检查注意力和量表理解，但不会通过“必须评 5 分”的规则机械排除参与者。本研究不包含高抑制档；其可见条纹和颗粒的可接受性留待未来工作评估（见 \S\ref{sec:discussion}）。

\textbf{样本量与分析计划}：计划至少招募 $N=24$ 名参与者，该样本量处于在 80\% 检验效能下检测 $d_z\approx0.6$--$0.8$ 配对效应量所需范围内，并且恰好覆盖两轮每 12 人一次的联合顺序分配。预注册的排除标准包括：调试/演示会话、未完整记录四次打字或六次评分、实测刷新率低于 200\,Hz、任一打字试验尝试字符少于 5 个、控制条件平均准确率低于 50\%、掩蔽试验未处于时序模式、任一试验的观测基本循环频率低于 50\,Hz，或以最短观看时长进行直线式评分。最后一项标准定义为：全部六行评分记录的观看时长均不超过 11~秒（在 10 秒提交门槛上容许 1 秒误差），且每行内部三个条目的分数完全相同。WPM、CPM、准确率、尝试字符数和首次按键延迟首先在每名参与者内部按条件取平均；速度指标必须与准确率联合解释，避免将低准确率下的快速输入误判为性能提升。当预先指定的 Shapiro 检验通过时，使用配对 $t$ 检验检验配对差异；否则使用 Wilcoxon 符号秩检验，并报告效应量和以参与者为单位的 bootstrap 95\% 置信区间。Likert 数据按维度使用 Friedman 检验，并以 Holm 校正的事后成对 Wilcoxon 检验进行分析。分析脚本直接从正式数据库 \texttt{study\_formal.db} 生成去标识化的参与者均值、JSON 审计报告和两个 \LaTeX\ 表格。

% NOTE: All "XX" marks and "--" table cells below are pending data placeholders.
% After the study, run webstudy/analyze_study.py; its analysis_output/typing_table.tex
% and ratings_table.tex rows map one-to-one onto Tables \ref{tab:study_typing} and
% \ref{tab:study_ratings}, and experiments/paper_figures/make_all.py regenerates
% Figs. \ref{fig:study_typing} and \ref{fig:study_ratings} from the same output.
\textbf{参与者}：共招募 XX 名志愿者；其中 XX 个会话满足全部预注册纳入标准并进入分析（排除 XX 个；按标准划分的审计计数随发布的分析输出一并提供）。纳入样本的可选人口统计信息如下：XX 名参与者中有 XX 名报告年龄（中位数 XX 岁，范围 XX--XX），XX 名报告性别（XX 名女性，XX 名男性）。所有正式会话均在 240\,Hz 实验室显示器上运行，实测刷新率至少为 200\,Hz。

\begin{table}[!t]
\caption{打字任务：控制条件与可读性优先掩蔽条件的被试内比较（$N=$ XX；每个条件取每名参与者两次试验的均值；$\Delta$ = 掩蔽 $-$ 控制，附参与者层面 bootstrap 95\% CI）。对每项指标，若预先指定的 Shapiro 检验通过，则采用配对 $t$ 检验（效应量为 Cohen's $d_z$）；否则采用 Wilcoxon 符号秩检验（效应量为秩二列相关 $r_{rb}$）；检验统计量在正文中报告}
\label{tab:study_typing}
\centering
\footnotesize
\setlength{\tabcolsep}{4pt}
\begin{tabular}{lccc}
\toprule
指标 & 控制 & 掩蔽 & $\Delta$ [95\% CI] \\
\midrule
准确率（\%） & -- & -- & -- [--, --] \\
WPM & -- & -- & -- [--, --] \\
CPM & -- & -- & -- [--, --] \\
尝试字符数 & -- & -- & -- [--, --] \\
首次按键延迟（ms） & -- & -- & -- [--, --] \\
\bottomrule
\end{tabular}
\end{table}

\begin{table}[!t]
\caption{六种显示条件下的主观评分（$N=$ XX；1--5 级 Likert 量表，越高越好；均值 $\pm$ 标准差）。稳定性为反向闪烁条目，舒适度为即时视觉舒适度条目（\S\ref{sec:user_study}）；Friedman 检验和 Holm 校正的成对 Wilcoxon 检验统计量在正文中报告}
\label{tab:study_ratings}
\centering
\footnotesize
\setlength{\tabcolsep}{4pt}
\begin{tabular}{lccc}
\toprule
条件 & 可读性 & 稳定性 & 舒适度 \\
\midrule
无掩蔽锚点 & -- & -- & -- \\
$n{=}2$ 掩蔽加噪声 & -- & -- & -- \\
$n{=}3$ 掩蔽加噪声 & -- & -- & -- \\
$n{=}4$ 掩蔽加噪声 & -- & -- & -- \\
$n{=}4$ 仅掩蔽 & -- & -- & -- \\
可读性优先完整配置（$n{=}4$） & -- & -- & -- \\
\bottomrule
\end{tabular}
\end{table}

\textbf{结果}：表~\ref{tab:study_typing} 和图~\ref{fig:study_typing} 汇总了打字任务。与控制条件相比，可读性优先完整配置使平均准确率变化 XX 个百分点（95\% CI [XX, XX]，$p=$ XX，效应量 XX），WPM 变化 XX（95\% CI [XX, XX]，$p=$ XX，效应量 XX），首次按键延迟变化 XX\,ms（95\% CI [XX, XX]，$p=$ XX，效应量 XX）；CPM 和尝试字符数对应的 $p$ 值分别为 XX 和 XX。表~\ref{tab:study_ratings} 和图~\ref{fig:study_ratings} 汇总了主观评分：可读性、稳定性和即时视觉舒适度的 Friedman 检验结果分别为 $\chi^2=$ XX（$p=$ XX）、$\chi^2=$ XX（$p=$ XX）和 $\chi^2=$ XX（$p=$ XX）；可读性优先完整配置与无掩蔽锚点在三个条目上的 Holm 校正成对比较 $p$ 值分别为 XX、XX 和 XX。

\begin{figure}[!t]
\centering
\includegraphics[width=\columnwidth]{figures/study_typing_paired.pdf}
\caption{控制条件与可读性优先掩蔽条件下的被试内打字表现：每条灰线代表一名参与者（每个条件下两次试验的均值）；标记表示组均值及参与者层面的 bootstrap 95\% 置信区间。各面板依次显示转录准确率、每分钟词数和首次按键延迟。}
\label{fig:study_typing}
\end{figure}

\begin{figure}[!t]
\centering
\includegraphics[width=\columnwidth]{figures/study_ratings_likert.pdf}
\caption{六种评分显示条件下的可读性、稳定性和即时视觉舒适度的 1--5 级 Likert 平均评分（误差线：标准差），条件从无掩蔽锚点排列至可读性优先完整配置。所有条目均为分值越高越好。}
\label{fig:study_ratings}
\end{figure}

}

% ----------------------------------------------------------
\subsection{补充检测与跟踪诊断}

目标检测与跟踪诊断在补充材料中报告，不纳入主要证据链。这些实验只使用一个真实拍摄位置、有限的 COCO 子集、MOT 静止帧拍摄和近似仿真指标，因此属于诊断性压力测试，不能扩展主要的传统 OCR 主张。主文省略归一化的跨任务汇总图，因为 OCR 字符恢复率、平均精度均值（mAP）、高阶跟踪准确率（HOTA）和身份识别 F1（IDF1）并非可直接比较的测量。

% ----------------------------------------------------------
\subsection{软件仿真补充实验}

本节使用软件仿真回答两个问题：（1）~经过应用渲染器后得到的单个受保护子帧，其机器可读性如何？（2）~在相机噪声为零且完全对齐时，恢复情况如何？0--2.6\% 的结果仅表征掩蔽后返回单个受保护子帧这一采集点。操作系统截图 API 可能在掩蔽前、合成后或跨已记录帧进行捕获；本研究未测试这些路径。截获原始帧的恶意软件，或对完整循环进行累积的录制，都可以绕过该方法。因此，我们不提出一般性的“抗截图”主张。

\subsubsection{合成单子帧 OCR}
在 120 个合成样本上使用 Tesseract~\cite{b_smith2007}、EasyOCR~\cite{b_easyocr2020} 和 Surya~\cite{b_surya2024} 后，三个引擎的字符恢复率从 94.0--97.0\% 降至 0--2.6\%（表~\ref{tab:ocr_corpus}）。该结果仅覆盖当前语料库、参数和引擎集合；分层细分可以检查类别趋势，但不足以证明对所有语言、布局或识别器均有效。

\begin{table}[!t]
\caption{合成单子帧 OCR：三引擎字符恢复率（120 个样本，$n=4$，$\varepsilon=8/255$）}
\label{tab:ocr_corpus}
\centering
\footnotesize
\begin{tabular}{lccc}
\toprule
OCR 引擎 & 原始帧 & 单子帧 & 降幅 \\
\midrule
Tesseract & 94.0\% & 0.0\% & 94.0\% \\
EasyOCR & 94.1\% & 0.0\% & 94.1\% \\
Surya & 97.0\% & 2.6\% & 94.4\% \\
\bottomrule
\end{tabular}
\end{table}

\begin{figure}[!t]
\centering
\includegraphics[width=\columnwidth]{figures/multiengine_ocr.pdf}
\caption{三个 OCR 引擎在原始帧和单个受保护子帧上的字符恢复率；误差线为以语料样本为单位重采样所得的 95\% bootstrap 区间。}
\label{fig:multiengine_ocr}
\end{figure}

图~\ref{fig:multiengine_ocr} 将这些结果可视化。所有引擎的单子帧恢复率均不超过 2.6\%，说明该发现并非由某一个识别器驱动。更广泛的推广仍需增加引擎和语料库。

\subsubsection{理论安全边界}
\label{sec:security_frontier}
当攻击者累积覆盖完整循环且完成配准的帧时，在式~(\ref{eq:noise_complement}) 的理想线性模型中，互补子帧可以重新求和。归档中逐样本最强的数字攻击恢复率为 95.2\%；纯完整循环时域平均的恢复率为 94.3\%。这是所评估线性分解的重建边界，而非正式的信息论定理。在 150 帧均值下，高抑制档的汇总恢复率仍为 47.9\%；排除欠曝光位置后为 53.9\%。字符恢复率不必与点亮像素比例成比例，因为识别结果还取决于内容和物理成像路径。

\begin{table}[!t]
\caption{归档的数字重建边界（Tesseract）}
\label{tab:digital_boundary}
\centering
\footnotesize
\setlength{\tabcolsep}{3pt}
\begin{tabular}{lrrrrr}
\toprule
设置 & $N$ & 字符恢复率 & 完全匹配率 & SSIM & $\Delta E_{00}$ \\
\midrule
基础：循环均值 & 120 & 94.3\% & 85.8\% & -- & -- \\
基础：逐样本最强 & 120 & 95.2\% & 91.7\% & -- & -- \\
可读性优先：循环均值 & 16 & 87.9\% & 50.0\% & 0.948 & 1.46 \\
高抑制：循环均值 & 16 & 56.6\% & 6.3\% & 0.764 & 4.87 \\
高抑制：反相槽 & 16 & 93.0\% & 81.2\% & -- & -- \\
\bottomrule
\end{tabular}
\vspace{1mm}

\raggedright\scriptsize 120 样本的基础攻击与 16 样本的档位消融是两项独立归档的数字评估。“逐样本最强”针对每个样本，从七种预先指定的重建中选择最佳结果。档位的 SSIM 和 $\Delta E_{00}$ 描述完整循环的重建质量，而非物理质量或用户感知质量。
\end{table}

主要证据链之外还保留了其他软件代理基线、参数扫描、检测和跟踪诊断。由于其物理条件、用户体验代价和指标与真实拍摄 OCR 的估计目标并不匹配，因此不将其作为主文证据。

\deferreduserstudy

%% ============================================================
%% VI. DISCUSSION
%% ============================================================
\section{讨论}
\label{sec:discussion}

\subsection{证据层级与主张范围}

最直接获得支持的结论是固定条件下的观测。在一条 S600 链路的共同标称设置下，所评估档位降低了三个开源 OCR 引擎在原始输入上的平均恢复率。固定的三引擎预处理网格能够恢复出更多文本：可读性优先档的网格最优选择达到 40.2\%，高抑制档达到 13.7\%，因此不能将原始输入视为攻击上界。相对于未保护输入，剩余的配对差值量化的是该网格内的降幅，而非对优化攻击者的抵抗能力。该结果既不构成曝光区间，也不构成一般性的防拍摄主张。语料库不平衡、仅使用一个网络摄像头、档位顺序固定，以及缺少亮度匹配的对照，使我们无法作出跨设备推断或机制隔离推断。VLM、积分和仿真结果用于界定边界，而非扩展主要主张。被动基线属于软件代理，不能证明本方法优于物理滤镜或既有时序显示系统。

\subsection{线性积分的数学张力}

该方案利用了人类时序感知与相机短曝光之间的差异。然而，其数学边界建立在理想的线性屏幕--相机辐亮度模型上，而不是将人类视觉视为线性求和算子。在该模型中，互补子帧可以重新求和以恢复原图（$\sum\mathbf{I}_k = \mathbf{I}$），因此能够累积一个完整配准循环的相机会击破这一线性分解。关于 VLM 对抗鲁棒性的研究~\cite{b_jiang2025} 进一步表明，大模型能够从受扰图像中提取部分语义信息；\S\ref{sec:vlm} 的探测实验验证了这一点。VLM 能够从单张短曝光照片中恢复大字号内容，并从 150 帧视频均值和长曝光视图中获得更强信号。在所有已评估重建中，归档数字评估的逐样本最强攻击达到 95.2\%；纯完整循环时域均值的恢复率达到 94.3\%（\S\ref{sec:security_frontier}）。

\subsection{物理层混杂因素与仿真--现实差距}

仿真--现实差距与若干物理层效应相容，但当前设计未能单独识别其中任何一种效应。

\textbf{方向 1：真实短曝光恢复率高于仿真（有利于攻击者）。}仿真单子帧 OCR 的字符恢复率为 0--2.6\%（表~\ref{tab:ocr_corpus}），而真实拍摄在相近掩蔽条件下的短曝光恢复率为 19.2\%（仅掩蔽档，表~\ref{tab:real_ocr_full_pool}）。两种物理效应可能对此有所贡献。第一，滚动快门行混合可能使不同行积分不同的子帧内容，但本研究未测量 S600 的读出时间。第二，LCD 响应转换可能混合相邻子帧，但未表征面板的灰阶转换行为。真实拍摄测量包含这两种效应，但本实验无法识别它们各自的贡献。因此，仿真是理想化的单子帧参照，而非物理下界。

\textbf{方向 2：真实积分恢复率低于仿真。}逐样本最强数字攻击的恢复率为 95.2\%（\S\ref{sec:security_frontier}），而真实 150 帧均值在可读性优先档和高抑制档上的恢复率分别为 71.1\% 和 47.9\%。共同设置下的数值分别为 79.7\% 和 53.9\%。ISP 处理、moir\'{e}（摩尔纹）、模糊和不完美配准都可能有所贡献，但归档无法识别它们是否存在及其幅度。专门设计的积分攻击因而可能缩小观测到的差距。

因此，仿真和物理实验回答的是不同问题。仿真评估理想化的单子帧变换；物理恢复率则测量包含未解决混杂因素在内的完整归档屏幕--相机流水线。

\subsection{探索性高抑制档的作用与代价}

探索性高抑制档通过条纹和字形扰动破坏严格完备性，并与较低的积分恢复率相关。其数字完整循环 SSIM 为 0.764，可读性优先档则为 0.948；$\Delta E_{00}$ 分别为 4.87 和 1.46（表~\ref{tab:digital_boundary}）。同一数字消融暴露出直接的单槽失效：经过对比度拉伸后，反相帧的平均字符恢复率为 93.0\%，完全匹配率为 81.2\%。这一观测到的 93.0\% 最佳数字恢复率是该档位最强的已观测数字攻击，因此 56.6\% 的时域平均结果并非最坏情形安全界限。在共同设置池中，高抑制档的短曝光字符恢复率为 5.6\%，150 帧均值下为 53.9\%。它没有达到描述性的 $<5\%$ 目标，且只能缓解积分泄漏。用户接受度尚未测试。

这一档位区分是解释结果的核心。可读性优先档是唯一与拟开展舒适度研究相联系的档位，但它在显式字段和积分攻击下更弱。高抑制档在若干视图中降低了观测恢复率，但可利用的反相槽和缺失的可用性证据使其不能支持一般性的安全主张。因此，本文支持的是一个可行性应用空间，而非部署建议。

\subsection{VLM 攻击者与对策方向}
\label{sec:vlm_defense}

VLM 的高恢复率集中在大字号短片段上（\S\ref{sec:vlm}），而时域平均能够提供更强信号。这一模式与保留的字形轮廓和语言先验补全相符，但本实验未单独识别其中任何一种机制。我们保留该结果，是因为它限制了方法的适用范围：降低传统 OCR 恢复率并不意味着能够抵抗 VLM。由此可以提出三个尚未验证的假设：字体大小自适应粒度、字形诱饵，以及字段级渲染策略。互补诱饵在理想的完整循环积分下仍会抵消，且这些设想均未经过物理验证。因此，实际目标应表述为提高攻击成本，而不是使可见信息无法恢复。

\subsection{持续视频攻击者}

所评估的视频攻击并未证明四帧即可绕过。拍摄流水线以标称 60\,fps 记录 150 个短曝光帧，报告的时域均值图像对完整连拍取平均，对应约 2.5\,s 的标称拍摄时长。由于持续平均会累积多个显示循环，该攻击不需要掩蔽种子、相位或精确的帧选择。其汇总恢复率在可读性优先档上为 71.1\%，在高抑制档上为 47.9\%；排除欠曝光位置后，分别升至 79.7\% 和 53.9\%。归档还包含从同一完整连拍中选取的聚合视图，但它们无法证明从独立采集的短窗口中也能恢复。时长--恢复率曲线和预先指定的短窗口分析尚未评估。因此，应区分单次拍摄适用范围与已测得的持续录制边界，同时不指定缺乏证据支持的最短绕过时长。

\subsection{部署路径}

真实部署需要驱动级帧缓冲注入、面板/背光协调和经过验证的时序。这些步骤会引入色彩管理、亮度饱和、丢帧和同步误差。在匹配短曝光分析中，可读性优先档的显式字段恢复率仍为 17.7\%，且在积分攻击下字符恢复率较高。因此，凭据保护不应仅依赖该档位。候选的字段级缓解措施需要分别进行可读性、曝光、视频和 VLM 评估。

\subsection{所需的验证性物理协议}

验证性采集应在每个内容--几何条件区块内随机化档位顺序，并记录墙钟时间批次标识符。实验应锁定并记录曝光模式、增益、白平衡、焦点、编解码器、固定的显示器亮度设置和环境照度。应使用光电二极管或高速传感器测量时隙时序和显示器物理亮度。每种时序档位都应配备在时间平均亮度和叠加外观上匹配的静态对照。这些对照可以将档位关联与时间漂移、占空比亮度、静态叠加、面板响应和相机处理区分开来。在获得此类数据之前，现有归档无法支持有关时序机制的因果主张。

\subsection{伦理声明}

本研究具有防御性质，旨在保护屏幕内容免遭未经授权的拍摄。攻击分析（\S\ref{sec:security_frontier}）旨在诚实评估安全边界，而非提供攻击指南。拟开展的用户研究属于最小风险的人体参与研究；目前尚未收集任何参与者数据。只有在具备权限的机构或院系主管部门出具书面批准、豁免或免审认定后，才会开始招募。协议规定了知情同意、光敏感筛查和自愿退出，详见 \S\ref{sec:user_study}。

\subsection{局限性}

\begin{enumerate}
\item 用户研究（\S\ref{sec:user_study}）目前是一项协议，而非已完成的证据：尚无参与者数据，因此可读性优先档的可用性仍未得到验证。即使在受控的 240\,Hz 实验室中完成研究，外部有效性仍需通过更广泛的参与者样本和显示条件加以检验。
\item 即时视觉舒适度是在短时暴露后采用单条目量表测得，不能推广至长期疲劳或临床视觉不适。
\item 该概念验证（PoC）仅位于应用层，不包含操作系统驱动程序或硬件背光。亮度、高动态范围（HDR）和色差结论来自数字模型，缺少光度学或面板响应测量。实际播放时隙时序未通过光电二极管、高速相机或示波器验证。由于实验缺少亮度匹配的静态对照和匹配的无噪声增强档位，因此无法将时序稀疏性、占空比亮度、静态叠加、面板响应、相机曝光和 ISP 处理作为机制区分开来。
\item 真实拍摄 OCR 仅使用一个 S600 UVC 网络摄像头。显示器亮度设置和采集相位保持固定，但未记录显示器物理亮度和环境照度。传感器增益、自动曝光语义、白平衡状态和焦点状态均未保留。档位顺序固定而非随机。因此，时间、批次和相机自动行为可能混杂档位或几何条件比较，但无法重建其方向和幅度。重复内容和拍摄重采样汇总不能替代跨设备复现。
\item 真实拍摄检测在单一位置使用了 150 张 COCO 图像；跟踪使用来自一个 MOT17 序列的 450 次静止帧拍摄，而非连续视频。仿真检测使用 8~张图像和近似指标。
\item VLM 评估覆盖一个相机、0$^{\circ}$ 正轴、一个档位和三个商业 API 模型（每个单元格 12 个视频片段）。模型可能随时间更新。公开文档内容可能通过训练数据记忆抬高恢复率。我们未评估多循环自适应重建或掩蔽前恶意软件截获。完整循环积分仍是所评估线性设计的重建边界。
\item 传统 OCR 评估覆盖 Surya、EasyOCR 和 Tesseract，但不包含商业 OCR API。预处理压力测试穷举了针对这些引擎预先声明的六输入网格，而非连续参数调优、学习式增强、图像复原或识别器特定优化。在所评估集合内，最优引擎与输入的联合选择有利于攻击者，但 Google Cloud Vision、Azure AI Vision、Baidu OCR 或未来识别器仍可能得到更高恢复率。
\item 基线比较使用未匹配的软件代理，而非经过物理测试的隐私滤镜或既有时序显示系统。这些比较用于提供观测权衡的背景，但不能证明本方法相对于这些替代方案具有比较优势。
\item 可读性优先档不适合作为敏感字段的唯一保护层：匹配短曝光条件下的字段微平均恢复率为 17.7\%，积分/VLM 攻击能恢复更多内容。仍需独立的访问控制。
\item 由 12 个内容项构成的语料库并非全英文/ASCII，而是混合语料：CET6 页面和两个混合语言片段包含中文字符。该语料库规模过小且经过选择性构造，无法支持文字系统层面的比较；专门的非拉丁文字评估仍属未来工作。
\end{enumerate}


%% ============================================================
%% VII. CONCLUSION
%% ============================================================
\section{结论}
\label{sec:conclusion}

我们在一个标称 UVC 控制设置下对显示处理进行了档位级测量。在 8 种共同设置几何条件下，每个档位使用相同的 288 个匹配单元（12 个内容项 $\times$ 8 种几何条件 $\times$ 3 次重复），未保护条件、可读性优先档和高抑制档的平均恢复率分别为 94.5\%、17.8\% 和 5.6\%。按内容项聚类的配对差异分别为 76.7 个百分点（[74.3, 79.2]）和 88.9 个百分点（[85.5, 91.9]）。这些结果支持三个开源 OCR 引擎在固定链路上的关联。它们既未确定曝光区间，也未确立时序机制的因果关系。可读性优先档仍有 17.7\% 的匹配字段微平均恢复率和 95.0\% 的 P99 字符恢复率，而高抑制档未达到其描述性的 $<5\%$ 目标。

绕过证据同样构成严格限制。150 帧时域均值在可读性优先档和高抑制档上的字符恢复率分别为 71.1\% 和 47.9\%；共同设置下的数值分别为 79.7\% 和 53.9\%。在完整的 9 几何条件敏感性池中，标称控制设置下的长曝光恢复率在可读性优先档上为 60.9\%，未保护条件为 47.3\%。正轴 VLM 探测也恢复了大量大字号内容。在完整固定网格上，三引擎网格最优选择在可读性优先档和高抑制档上分别达到 40.2\% 和 13.7\%。不受约束或学习式预处理以及商业 OCR 均不在当前证据所支持的条件范围内。

因此，本研究贡献的是一项有边界的物理测量及明确的失效约束，而非可部署的防相机方案。下一步最关键的实验是采用匹配静态对照的随机化采集、曝光--恢复率曲线、可用性测试和跨设备复现。

%% ============================================================
%% DATA AND CODE AVAILABILITY
%% ============================================================
\section*{数据与代码可用性}

软件、实验配置、去标识化的逐样本 OCR 结果、VLM 记录和分析脚本维护于 \url{https://github.com/docandyc/privacy-display}。该仓库目前仍可变更，不作为不可变的投稿产物提供。为支持复现，定稿版本将提供冻结的发布标签或归档 DOI。VLM 记录包含端点标识符、时间戳、提示词、解码参数、输入哈希或路径、解析状态和原始输出。归档标识符 \texttt{capture\_hardened} 和 \texttt{vlm} 均表示高抑制档。若用户研究得以开展，可识别身份的注册数据将不予公开；仅公开去标识化的结果和审计信息。

%% ============================================================
%% ACKNOWLEDGMENT
%% ============================================================
\section*{致谢}
\TODO{致谢待填写。}

OpenAI Codex~\cite{b_openai_codex2026} 协助开发分析脚本、执行一致性检查，并对摘要、引言、实验评估、讨论和结论进行英文修订。作者确定科学决策，依据归档数据核验生成的分析，审阅所有保留的文本与代码，并对本文承担责任。


%% ============================================================
%% REFERENCES
%% ============================================================
\bibliographystyle{IEEEtran}
\bibliography{refs}
\EOD

\vfill

\end{document}
